\chapter{Performance Analysis and Results}

\section{Experimental Setup}

All performance benchmarks were conducted on the following system configuration:
\begin{itemize}
\item Operating System: Windows 11
\item Python Version: 3.x with NumPy for numerical operations
\item CPU: Modern x86-64 processor
\end{itemize}

The performance comparison script executes multiple iterations of key generation and key exchange/encapsulation operations for both ECC and Kyber implementations (512, 768, and 1024-bit security levels). Each test was run multiple times to obtain statistically meaningful results.

\section{Performance Metrics}

\subsection{Execution Time Comparison}

Table \ref{tab:exec-time} compares the execution times between ECC and Kyber implementations:

\begin{table}[!htbp]
\centering
\caption{\label{tab:exec-time}Mean Execution Time (milliseconds)}
\begin{tabular}{l|r|r}
\hline
\textbf{Algorithm} & \textbf{Key Generation} & \textbf{Exchange/Encapsulation} \\ \hline
ECC (secp256k1) & 97.76 & 95.62 \\
Kyber-512 & 70.54 & 107.15 \\
Kyber-768 & 128.44 & 172.85 \\
Kyber-1024 & 199.47 & 252.12 \\ \hline
\end{tabular}
\end{table}

Kyber-512 performs comparably to ECC in key generation (70.54 ms vs 97.76 ms) but shows slower encapsulation times (107.15 ms vs 95.62 ms). This is expected given the pure Python implementation and the computational complexity of lattice-based operations compared to elliptic curve arithmetic.

\subsection{Performance Relative to ECC Baseline}

Table \ref{tab:slowdown} presents the performance of each algorithm relative to ECC:

\begin{table}[!htbp]
\centering
\caption{\label{tab:slowdown}Performance Factor Relative to ECC}
\begin{tabular}{l|r|r}
\hline
\textbf{Algorithm} & \textbf{Key Generation} & \textbf{Encapsulation} \\ \hline
ECC (baseline) & 1.00x & 1.00x \\
Kyber-512 & 0.72x & 1.12x \\
Kyber-768 & 1.31x & 1.81x \\
Kyber-1024 & 2.04x & 2.64x \\ \hline
\end{tabular}
\end{table}

Values below 1.0x indicate faster performance than ECC. Kyber-512 key generation outperforms ECC (0.72x), while encapsulation lags slightly behind (1.12x). Higher security levels show proportionally increased execution times due to larger matrix dimensions and polynomial operations.

\section{Key Size Comparison}

\begin{table}[!htbp]
\centering
\caption{\label{tab:keysize}Key and Ciphertext Sizes (bytes)}
\begin{tabular}{l|r|r|r}
\hline
\textbf{Algorithm} & \textbf{Public Key} & \textbf{Ciphertext} & \textbf{Total Bandwidth} \\ \hline
ECC (secp256k1) & 33 & 33 & 66 \\
Kyber-512 & 800 & 768 & 1,568 \\
Kyber-768 & 1,184 & 1,088 & 2,272 \\
Kyber-1024 & 1,568 & 1,568 & 3,136 \\ \hline
\end{tabular}
\end{table}

While Kyber requires larger keys and ciphertexts compared to ECC, the sizes remain practical for modern network communications. The bandwidth overhead factor ranges from 23.8x for Kyber-512 to 47.5x for Kyber-1024 relative to ECC.

\subsection{Memory Usage}

Table \ref{tab:memory} shows peak memory consumption during cryptographic operations:

\begin{table}[!htbp]
\centering
\caption{\label{tab:memory}Peak Memory Usage (MB)}
\begin{tabular}{l|r|r}
\hline
\textbf{Algorithm} & \textbf{Key Generation} & \textbf{Encapsulation} \\ \hline
ECC & 0.00 & 0.00 \\
Kyber-512 & 0.11 & 0.13 \\
Kyber-768 & 0.18 & 0.20 \\
Kyber-1024 & 0.28 & 0.29 \\ \hline
\end{tabular}
\end{table}

Memory usage remains modest across all implementations, with even Kyber-1024 requiring approximately 0.3 MB for operations—well within the capabilities of resource-constrained devices.

\section{Analysis and Discussion}

\subsection{Implementation Characteristics}

The benchmark results reflect the pure Python implementation using NumPy for array operations. Both ECC and Kyber implementations prioritize correctness and code clarity over raw performance, making them suitable for educational purposes and prototype systems.

Kyber-512 shows competitive performance with ECC despite using more complex lattice operations. Key generation is actually faster than ECC (0.72x), though encapsulation is slightly slower (1.12x). This demonstrates that lattice-based cryptography can achieve practical performance even without low-level optimizations.

Key implementation aspects:
\begin{itemize}
\item Number Theoretic Transform (NTT) for efficient polynomial multiplication
\item Modular arithmetic using Python's native integer operations
\item NumPy arrays for coefficient storage and basic operations
\item Standard cryptographic hash functions (SHA3, SHAKE) from hashlib
\end{itemize}

The polynomial operations in Kyber are inherently parallelizable. Matrix-vector multiplications operate on independent coefficients, and NTT computations follow regular patterns. This architectural characteristic means the algorithm would benefit significantly from optimization efforts in future work.

\subsection{Performance Scaling Across Security Levels}

As security levels increase, performance scales proportionally with the increased matrix and polynomial dimensions:

\begin{itemize}
\item \textbf{Kyber-512:} Fastest option, competitive with ECC (quantum-resistant equivalent to AES-128 and secp256k1)
\item \textbf{Kyber-768:} Balanced performance, 1.31x-1.81x relative to ECC (quantum-resistant equivalent to AES-192 and secp384r1)
\item \textbf{Kyber-1024:} Highest security, 2.04x-2.64x relative to ECC (quantum-resistant equivalent to AES-256 and secp521r1)
\end{itemize}

For most applications, Kyber-512 provides adequate security (128-bit, matching widely-deployed secp256k1) with reasonable performance characteristics in pure Python implementation.

\subsection{Portability and Platform Independence}

The pure Python implementation ensures broad portability across platforms. The implementation runs on any system with Python 3.x and NumPy, including:
\begin{itemize}
\item x86-64 desktop and server systems (Windows, Linux, macOS)
\item ARM devices including mobile phones and single-board computers
\item Cross-platform environments without specific hardware requirements
\end{itemize}

This platform independence makes the implementation particularly suitable for:
\begin{itemize}
\item Educational purposes and algorithm understanding
\item Rapid prototyping and testing
\item Cross-platform applications requiring consistent behavior
\item Systems where native compilation is impractical
\end{itemize}

The trade-off is performance: production systems would benefit from optimized implementations in C/C++ or Rust with platform-specific optimizations. However, for many use cases, the current performance is acceptable, especially for infrequent operations like TLS handshakes or key establishment.

\subsection{Bandwidth Trade-offs}

The primary disadvantage of Kyber is bandwidth overhead. Table \ref{tab:bandwidth} summarizes the bandwidth increase:

\begin{table}[!htbp]
\centering
\caption{\label{tab:bandwidth}Bandwidth Overhead Relative to ECC}
\begin{tabular}{l|r|r}
\hline
\textbf{Algorithm} & \textbf{Total Bytes} & \textbf{Overhead Factor} \\ \hline
ECC & 66 & 1.0x \\
Kyber-512 & 1,568 & 23.8x \\
Kyber-768 & 2,272 & 34.4x \\
Kyber-1024 & 3,136 & 47.5x \\ \hline
\end{tabular}
\end{table}

However, in modern networks with high bandwidth availability, this overhead is manageable. For a single key exchange, even Kyber-1024's 3,136 bytes represents less than 0.003 seconds of transmission time on a 10 Mbps connection.

\section{Deployment Considerations}

The transition from classical algorithms to Post-Quantum Cryptography represents a fundamental shift in cryptographic infrastructure. Based on our experimental results, this section analyzes the expected operational changes and necessary adjustments for deploying Kyber in production environments.

\subsection{Latency and Throughput Trade-offs}

The data shows that post-quantum cryptography can achieve practical performance even in pure Python. While Kyber introduces a significant increase in bandwidth usage (approximately 23.8x overhead for Kyber-512 relative to ECC), the computational performance is competitive.

Kyber-512 shows balanced performance characteristics: faster key generation (0.72x) but slightly slower encapsulation (1.12x) compared to ECC. For high-bandwidth environments, the computational characteristics combined with the smaller keygen time can still result in acceptable overall handshake performance. However, in bandwidth-constrained environments (e.g., IoT/LoRaWAN), the transmission time for larger keys may become the bottleneck.

Systems migrating to Kyber must account for Maximum Transmission Unit (MTU) fragmentation, as Kyber keys and ciphertexts exceed the standard Ethernet MTU of 1500 bytes, potentially requiring packet fragmentation handling at the network layer \cite{Paquin2020}.

\subsection{Hybrid Implementation Strategies}

Immediate replacement of classical algorithms with Kyber is generally not recommended due to the relative newness of lattice-based cryptography compared to ECC. A hybrid deployment is the standard migration path endorsed by security bodies including NIST and BSI.

In a hybrid scheme, the final shared secret is derived from both an ECC key exchange and a Kyber encapsulation:
\[ K_{final} = \text{KDF}(K_{ECC} || K_{Kyber}) \]
This approach ensures that security remains at least as strong as the classical baseline even if a vulnerability is later discovered in the lattice construction. While our benchmarks treated Kyber as a standalone replacement, real-world deployment would likely incorporate both schemes. A hybrid handshake would require both ECC (98 ms KeyGen + 96 ms Exchange = 194 ms) and Kyber-512 (71 ms KeyGen + 107 ms Encapsulation = 178 ms) for a total of approximately 372 ms. Compared to ECC-only deployment, this represents a 92\% increase in handshake time while providing quantum resistance. This overhead is manageable for most applications, making hybrid adoption practically viable \cite{Bernstein2017,Paquin2020}.

\subsection{Operational Stability and Statistical Consistency}

To ensure the reliability of these performance claims, our experimental methodology prioritized statistical significance over single-run metrics. The reported results represent the mean of $N=1000$ iterations for each security level, minimizing the impact of background factors.

The multiple iterations reveal consistent timing characteristics for both algorithms. Kyber's lattice operations (matrix-vector multiplication and NTT) follow regular computational patterns without data-dependent branching, while ECC point operations involve modular arithmetic with variable execution paths. For real-time systems requiring predictable execution times, both algorithms demonstrate acceptable consistency when implemented with constant-time techniques.